\chapter{Architettura esagonale del sistema}
%\epigraph{Citazione}{Autore della citazione}
\label{appendix:hexagonal-architecture}

\section{Introduzione}

L’architettura di ogni microservizio del \textit{backend} è stata progettata seguendo il paradigma
dell’architettura esagonale (\textit{Hexagonal Architecture}, o \textit{Ports and Adapters}).
Questo approccio consente di separare nettamente la logica di dominio
dalle dipendenze esterne, come database, \gls{apig}, \gls{frameworkg} web e servizi esterni.

L’obiettivo principale di questa scelta architetturale è garantire il rispetto
dei principi SOLID e, più in generale, promuovere una progettazione modulare
e manutenibile del sistema. In particolare, l’architettura adottata consente:
\begin{itemize}
    \item \textbf{Isolamento della logica di business}, che rimane indipendente
    da framework, database e tecnologie esterne;

    \item \textbf{Rispetto del principio di Single Responsibility (SRP)},
    grazie alla separazione tra dominio, applicazione e infrastruttura;

    \item \textbf{Rispetto del principio di Dependency Inversion (DIP)},
    ottenuto tramite l’uso di \textit{ports} e \textit{adapters}, che permettono
    al dominio di dipendere da astrazioni e non da implementazioni concrete;

    \item \textbf{Testabilità elevata}, grazie alla possibilità di sostituire
    facilmente le dipendenze esterne con \textit{mock} o implementazioni fittizie;

    \item \textbf{Manutenibilità ed estensibilità}, facilitando l’evoluzione del
    sistema senza impatti significativi sulle componenti esistenti;

    \item \textbf{Disaccoppiamento tra dominio e infrastruttura}, che consente
    di modificare dettagli tecnici (database, \gls{llmg}, servizi esterni) senza
    alterare la logica interna al sistema.
\end{itemize}
Nel sistema sviluppato, tale paradigma è stato applicato in modo coerente
a tutti i microservizi principali: \textit{auth-service}, \textit{user-service},
\textit{search-service}, \textit{history-service} e \textit{mcpserver-service}.

\section{Struttura generale dei microservizi}
Ogni microservizio del sistema segue una struttura standardizzata, 
organizzata nei seguenti livelli:
\begin{itemize}
    \item \textbf{Domain}: contiene le entità, i \textit{value object}, i servizi di dominio e le eccezioni;
    \item \textbf{Application}: definisce i casi d’uso e le interfacce (ports);
    \item \textbf{Adapters}: implementa l’interazione con l’esterno (\textit{inbound} e \textit{outbound});
    \item \textbf{Infrastructure}: gestisce dettagli tecnici come persistenza, DTO (\textit{Data Transfer Object}) e integrazioni.
\end{itemize}


\section{Architettura del progetto}
La struttura del progetto segue il seguente schema ricorrente:
\begin{verbatim}
src/
 ├── adapters
 │    ├── inbound
 │    └── outbound
 ├── application
 │    ├── ports
 │    │    ├── inbound
 │    │    └── outbound
 │    └── repositories
 ├── domain
 │    ├── entities
 │    ├── services
 │    ├── value_objects
 │    └── exceptions
 ├── infrastructure
 │    ├── persistence
 │    └── dtos
 └── main.py
\end{verbatim}
Per comprendere l'efficacia del paradigma \textit{Ports and Adapters} 
all'interno del sistema è necessario analizzare la responsabilità di ciascun livello 
e come questi interagiscono tra loro.

\subsection{Il Core: Domain e Application}
Il centro dell'esagono è costituito dal \textbf{Domain} e dal livello di \textbf{Application}. 
Questo nucleo non ha alcuna conoscenza del mondo esterno: non sa se i dati arrivino da 
una chiamata \gls{httpg}, da una riga di comando o da un test unitario. 
Tale livello si compone dei seguenti elementi:
\begin{itemize}
    \item \textbf{Domain Entities e Value Objects}: rappresentano i concetti di business puri. 
    Ad esempio, in History Service, \texttt{ConversationEntity} e \texttt{MessageEntity} 
    contengono la struttura per la definizione della conversazione.
    \begin{lstlisting}[language=Python, caption={Definizione delle classi ConversationEntity e MessageEntity}, label={lst:domain-entities}]
@dataclass(frozen=True)
class ConversationEntity:
    id: UUID
    user_id: UUID
    created_at: datetime
    messages: list[MessageEntity]

@dataclass(frozen=True)
class MessageEntity:
    id: UUID
    conversation_id: UUID
    role: str
    content: str
    type: str
    format: str
    created_at: datetime
\end{lstlisting}
    \item \textbf{Domain Services}: incapsulano logiche di business complesse che coinvolgono più entità. 
    Un esempio è il \texttt{CrawlToolService} in MCP Server Service, che coordina l'algoritmo
    del \gls{toolg} di \gls{crawlingg} senza legarsi a client \gls{httpg} o librerie di 
    \textit{scraping} specifiche.
    \begin{lstlisting}[language=Python, caption={Implementazione di CrawlToolService}, label={lst:auth_controller}]
@dataclass(frozen=True)
class CrawlToolService(CrawlToolUseCase):
    cachePort: CachePort
    fetchPagePort: FetchPagePort
    summarizePort: SummarizePort
    checkAnswerabilityPort: CheckAnswerabilityPort
    chooseLinksPort: ChooseLinksPort
    formatAnswerPort: FormatAnswerPort
    notifyProgressPort: NotifyProgressPort
    manageChunksPort: ManageChunksPort

    def execute(self, param):
        # Logica di crawling
    \end{lstlisting}
    \item \textbf{Inbound Ports (Use Cases)}: sono interfacce che definiscono ciò che l'applicazione 
    può fare. 
    Nel microservizio di autenticazione, per esempio, \texttt{AuthUseCase} espone il contratto per le operazioni 
    di login o di rinnovo del \textit{token}.
    \begin{lstlisting}[language=Python, caption={Definizione di AuthUseCase}, label={lst:auth_controller}]
class AuthUseCase(ABC):
    @abstractmethod
    def login(self, username: str, password: str) -> TokenPair : ...

    @abstractmethod
    def refresh(self, token: str) -> TokenPair : ...

    @abstractmethod
    def logout(self, token: str) -> None : ...

    @abstractmethod
    def getJwks(self) -> dict : ...
    \end{lstlisting}
    \item \textbf{Outbound Ports}: sono interfacce che definiscono i bisogni infrastrutturali del dominio. 
    Ad esempio, il dominio di MCP Server Service ha bisogno di archiviare i dati estratti o di 
    notificare l'avanzamento dei passaggi all'utente; per farlo, definisce astrazioni come 
    \texttt{CachePort}.
    \begin{lstlisting}[language=Python, caption={Definizione di CachePort}, label={lst:auth_controller}]
class CachePort(ABC):
    @abstractmethod
    def findCachedRequest(self, param: CrawlParam) -> ToolResult | None : ...

    @abstractmethod
    def storeRequest(self, param: CrawlParam, result: str) -> None : ...
    \end{lstlisting}
\end{itemize}

\subsection{I Punti di Ingresso: Inbound Adapters}
Gli \textit{Inbound Adapters} (o \textit{Driving Adapters}) guidano l'applicazione catturando gli 
stimoli esterni e traducendoli in comandi per i casi d'uso interni. 
Nel progetto, questi adattatori sono rappresentati dai \textbf{Controller} web FastAPI. 

Prendendo come esempio la classe \texttt{AuthController} all'interno di Auth Service:
\begin{enumerate}
    \item Il controller espone l'\textit{endpoint} \gls{httpg} \texttt{POST /login};
    \item Riceve la richiesta formattata secondo il \texttt{LoginDto};
    \item Valida l'input sintattico e mappa il DTO in un oggetto di dominio;
    \item Invoca l'\textit{inbound port} \texttt{AuthUseCase} per eseguire la logica di autenticazione.
\end{enumerate}
\begin{lstlisting}[language=Python, caption={Definizione endpoint "/login" in AuthController}, label={lst:auth_controller}]
class AuthController:
    def __init__(self, authUseCase: AuthUseCase):
        self.authUseCase = authUseCase
        self.router = APIRouter()
        self._register_routes()

    def _register_routes(self):
        @self.router.post("/login")
        def login(dto: LoginDto):
            try:
                return self.authUseCase.login(dto.username, dto.password)
            except (UnauthorizedException, UserNotFoundException):
                raise HTTPException(status_code=403, detail="Forbidden")
            except Exception as e:
                raise HTTPException(status_code=500, detail=str(e))

        [...]
\end{lstlisting}

\subsection{Le Connessioni Esterne: Outbound Adapters e Persistence}
Gli \textit{Outbound Adapters} (o \textit{Driven Adapters}) implementano le interfacce 
definite dalle porte d'uscita, gestendo l'effettiva interazione con i dettagli 
tecnologici. 
Grazie al meccanismo della \textit{Dependency Inversion}, le classi infrastrutturali 
ereditano dalle porte dell'applicazione.
Un esempio lampante di questo disaccoppiamento è visibile in MCP Server Service:
\begin{itemize}
    \item L'applicazione definisce la porta \texttt{SummarizePort} per richiedere il 
    riassunto di un testo;
    \item All'esterno dell'esagono, l'adattatore \texttt{SummarizeAdapter} 
    implementa tale interfaccia interfacciandosi con le \gls{apig} 
    di OpenAI.
\end{itemize}
Lo stesso principio regola il livello di persistenza dei dati: i moduli 
\texttt{Repository} definiscono le interfacce e i contratti astratti di accesso ai dati, 
la cui implementazione concreta è demandata alle rispettive classi \texttt{RepositoryImpl} 
all'interno dello strato infrastrutturale. 
Grazie a questo totale disaccoppiamento, 
qualora in futuro l'azienda decidesse di sostituire l'attuale database relazionale, 
l'intervento di \textit{refactoring} si limiterebbe alla sola riscrittura della classe di 
implementazione, lasciando il nucleo completamente inalterato.

Il medesimo vantaggio architetturale si riscontra nella gestione dei modelli di linguaggio: 
un'eventuale transizione verso un \gls{llmg} differente da quello attuale comporterebbe 
esclusivamente l'aggiornamento o la sostituzione dello specifico \textit{outbound adapter}, 
senza alcun impatto sulla logica di calcolo e sul comportamento globale del sistema.